\chapter{Méthodologie et outils utilisés}

\section{Approche de travail}

\noindent
Le développement de \projectName{} a été réalisé de manière progressive, en partant des besoins essentiels du projet puis en améliorant l'expérience utilisateur et la stabilité de l'application. Cette approche a permis de construire d'abord les fonctionnalités principales, puis de corriger les problèmes de logique, de navigation, de cohérence visuelle et de qualité du code.

\noindent
Le travail a été organisé autour de cycles courts :

\begin{itemize}
    \item analyse du besoin et identification des rôles ;
    \item conception des données et des routes API ;
    \item développement des interfaces principales ;
    \item intégration entre frontend et backend ;
    \item tests des parcours utilisateur ;
    \item corrections et stabilisation pour la démonstration.
\end{itemize}

\section{Outils de développement}

\noindent
Plusieurs outils ont été utilisés pour réaliser le projet :

\begin{table}[h!]
\centering
\renewcommand{\arraystretch}{1.4}
\begin{tabular}{|p{4cm}|p{9cm}|}
\hline
\textbf{Outil} & \textbf{Utilisation} \\
\hline
Visual Studio Code & Édition du code frontend, backend et du rapport. \\
\hline
Git et GitHub & Versionnement du projet et sauvegarde du code source. \\
\hline
Laragon & Environnement local pour exécuter Laravel et MySQL. \\
\hline
Postman / navigateur & Vérification manuelle des API et des parcours. \\
\hline
Composer & Gestion des dépendances PHP Laravel. \\
\hline
npm & Gestion des dépendances frontend React. \\
\hline
\end{tabular}
\caption{Outils utilisés durant le développement}
\label{tab:outils_dev}
\end{table}

\section{Technologies frontend}

\subsection{React et Vite}

\noindent
Le frontend de \projectName{} est développé avec \textbf{React}. Cette bibliothèque permet de construire des interfaces dynamiques à partir de composants réutilisables. \textbf{Vite} a été utilisé comme outil de développement et de build, notamment pour sa rapidité et sa configuration légère.

\subsection{Tailwind CSS}

\noindent
\textbf{Tailwind CSS} a été utilisé pour construire rapidement une interface responsive et cohérente. Il a permis de mettre en place l'identité visuelle de SmartJobs : thème sombre premium, accent orange, cartes avec effet glassmorphism et adaptation au mode clair.

\subsection{Framer Motion}

\noindent
\textbf{Framer Motion} a été utilisé pour améliorer l'expérience utilisateur grâce à des animations légères, notamment sur les transitions, les cartes et certains composants interactifs.

\section{Technologies backend}

\subsection{Laravel 11}

\noindent
Le backend est développé avec \textbf{Laravel 11}. Il fournit une API REST structurée permettant la gestion de l'authentification, des offres, des candidatures, des quiz, des notifications et des actions d'administration.

\subsection{Laravel Sanctum}

\noindent
\textbf{Laravel Sanctum} est utilisé pour sécuriser l'accès aux endpoints protégés. Les utilisateurs authentifiés accèdent aux fonctionnalités selon leur rôle : candidat, recruteur ou administrateur.

\subsection{MySQL}

\noindent
La base de données \textbf{MySQL} stocke les utilisateurs, profils, offres, candidatures, quiz, notifications, messages et paiements. Les migrations Laravel permettent de gérer l'évolution du schéma de manière contrôlée.

\section{Qualité et validation}

\noindent
La qualité du projet a été assurée à travers plusieurs vérifications :

\begin{itemize}
    \item compilation frontend avec \texttt{npm run build} ;
    \item analyse statique avec \texttt{npm run lint} ;
    \item tests backend Laravel avec \texttt{php artisan test} ;
    \item tests manuels des parcours candidat, recruteur et administrateur ;
    \item vérification des règles métier critiques comme l'unicité des candidatures et l'accès aux quiz.
\end{itemize}

\section{Organisation du code}

\noindent
Le projet est organisé en deux parties :

\begin{itemize}
    \item \textbf{frontend/} : application React, pages, composants, styles, appels API ;
    \item \textbf{backend/} : API Laravel, contrôleurs, modèles, migrations, tests et routes.
\end{itemize}

\noindent
Cette séparation facilite la maintenance, la lisibilité et l'évolution future de la plateforme.
